\subsection{Variables}

\subsubsection{Regla: camelCase y nombres descriptivos}

\textbf{Regla:} Variables locales y parámetros en \texttt{camelCase}. Se permite \texttt{i} únicamente como índice de un ciclo corto, simple y no anidado. Se permiten \texttt{j} y \texttt{k} solo como segundo y tercer índice de una operación posicional o matemática breve, sin significado de dominio. Cuando el contador representa filas, columnas, jugadores, rondas u otro concepto del juego, o cuando el cuerpo del ciclo no es inmediato, se usa un nombre descriptivo como \texttt{rowIndex}, \texttt{columnIndex} o \texttt{playerIndex} (Microsoft, 2025c).

\textbf{Código correcto:}
\begin{lstlisting}[style=csharpstyle]
public void RevealTiles(IReadOnlyList<Tile> tiles)
{
    for (int i = 0; i < tiles.Count; i++)
    {
        tiles[i].Reveal();
    }
}

public void ResetBoard()
{
    for (int rowIndex = 0; rowIndex < BoardSize; rowIndex++)
    {
        for (int columnIndex = 0; columnIndex < BoardSize; columnIndex++)
        {
            _tiles[rowIndex, columnIndex].Reset();
        }
    }
}
\end{lstlisting}

\textbf{Código incorrecto:}
\begin{lstlisting}[style=csharpstyle]
public void ResetBoard()
{
    for (int i = 0; i < BoardSize; i++)
    {
        for (int j = 0; j < BoardSize; j++)
        {
            _tiles[i, j].Reset();
        }
    }
}
\end{lstlisting}

\subsubsection{Regla: nombres de colecciones en plural}

\textbf{Regla:} Toda colección (arreglo, lista, diccionario o cualquier \texttt{IEnumerable<T>}) se nombra con un sustantivo en plural (Cwalina \& Abrams, 2008).

\textbf{Código correcto:}
\begin{lstlisting}[style=csharpstyle]
public sealed class GameBoard
{
    private readonly List<Player> _players;

    public GameBoard(List<Player> players)
    {
        _players = players;
    }
}
\end{lstlisting}

\textbf{Código incorrecto:}
\begin{lstlisting}[style=csharpstyle]
public sealed class GameBoard
{
    private readonly List<Player> _playerList;

    public GameBoard(List<Player> playerList)
    {
        _playerList = playerList;
    }
}
\end{lstlisting}

\subsubsection{Regla: una declaración por línea}

\textbf{Regla:} No se declara más de una variable en la misma sentencia, aunque compartan tipo (McConnell, 2004).

\textbf{Código correcto:}
\begin{lstlisting}[style=csharpstyle]
int currentRow = startingRow;
int currentColumn = startingColumn;
\end{lstlisting}

\textbf{Código incorrecto:}
\begin{lstlisting}[style=csharpstyle]
int currentRow = startingRow, currentColumn = startingColumn;
\end{lstlisting}

\subsubsection{Regla: prefijos interrogativos para variables booleanas}

\textbf{Regla:} Toda variable \texttt{bool} se nombra como una pregunta cerrada, con prefijo \texttt{is}, \texttt{has}, \texttt{can} o \texttt{should} (Martin, 2008).

\textbf{Código correcto:}
\begin{lstlisting}[style=csharpstyle]
bool isPlayerTurnActive = true;
bool hasReachedFinalTile = false;
\end{lstlisting}

\textbf{Código incorrecto:}
\begin{lstlisting}[style=csharpstyle]
bool playerTurnActive = true;
bool finalTileReached = false;
\end{lstlisting}
